\subsection{局部紧空间上在无穷远处为 0 的连续函数}

在本小节中，X 是一个局部紧空间。我们用 \(\mathcal{C}_{0}(X)\) 表示 X 上在无穷远处为 0 的复值连续函数所成的交换 Banach 代数，带上范数 \(\|f\|=\sup_{x\in X}|f(x)|\) (I, p.~17 的例 3)。

命题 1. --- 对 X 的每个闭子集 \(\Phi\)，用 \(I_{\Phi}\) 表示由 \(f \in \mathcal{C}_0(X)\)（在 \(\Phi\) 上为 0）所成的集合。则 \(\Phi \mapsto I_{\Phi}\) 是从 X 的闭子集所成的集合到 \(\mathcal{C}_0(X)\) 的闭理想所成的集合上的一个双射。

集合 \(\mathrm{I}_{\Phi}\) 是 \(\mathcal{C}_0(\mathrm{X})\) 的一个闭理想。

设 \(\Phi \neq \Phi'\) 是 X 的两个闭子集。如有必要，交换 \(\Phi\) 与 \(\Phi'\)，我们就可以假定存在 \(x \in \Phi'\) 使 \(x \notin \Phi\)；于是存在函数 \(f \in \mathcal{C}_0(X)\)，它在 \(\Phi\) 上为 0 且在 \(x\) 处非零 (TG, IX, p.~43, 命题 1)。我们有 \(f \in I_{\Phi}\) 且 \(f \notin I_{\Phi'}\)，因而映射 \(\Phi \mapsto I_{\Phi}\) 是单射。

设 I 是 \(\mathcal{C}_{0}(\mathrm{X})\) 的一个闭理想。设 \(\Phi\) 是由 \(x \in X\)（使 \(f(x) = 0\) 对一切 \(f \in I\) 成立）所成的集合；它是 X 的一个闭子集，且我们有 \(I \subset I_{\Phi}\)。我们来证明 \(I_{\Phi} \subset I\)，这将蕴涵 \(I = I_{\Phi}\) 并完成命题的证明。

设 \(f \in \mathrm{I}_{\Phi}\)。对每个实数 \(\varepsilon > 0\)，用 \(\mathrm{C}_{\varepsilon}\) 表示由 \(x \in \mathrm{X}\)（满足 \(|f(x)| \geqslant \varepsilon\)）所成的集合。由于 \(f\) 在无穷远处趋于 0，集合 \(\mathrm{C}_{\varepsilon}\) 是紧的。设 \(x \in \mathrm{C}_{\varepsilon}\)；由于 \(f(x) \neq 0\) 且 \(f \in \mathrm{I}_{\Phi}\)，有 \(x \notin \Phi\)；于是由 \(\Phi\) 的定义，存在函数 \(\varphi_x \in \mathrm{I}\) 使 \(|\varphi_x(x)| > 1\)，从而使得 \(|\varphi_x(y)| > 1\) 对一切 \(y\)（属于 \(\mathrm{V}_x\)，即 \(x\) 的一个邻域）成立。开集 \(\mathrm{V}_x \cap \mathrm{C}_{\varepsilon}\) 覆盖 \(\mathrm{C}_{\varepsilon}\)。由于集合 \(\mathrm{C}_{\varepsilon}\) 是紧的，存在有限子集 \(\mathrm{T}_{\varepsilon} \subset \mathrm{X}\)，使得

\[
\mathrm{C} _ {\varepsilon} \subset \bigcup_ {x \in \mathrm{T} _ {\varepsilon}} \mathrm{V} _ {x}.
\]

于是，元素

\[
g _ {\varepsilon} = \frac {1}{\varepsilon} \sum_ {x \in \mathrm{T} _ {\varepsilon}} \varphi_ {x} \overline {{\varphi_ {x}}} \geqslant 0
\]

（\(\mathcal{C}_0(\mathrm{X})\) 中的）属于 I，且我们有 \(g_{\varepsilon} \geqslant \varepsilon^{-1}\)（在 \(\mathrm{C}_{\varepsilon}\) 上）。函数

\[
f _ {\varepsilon} = \frac {f g _ {\varepsilon}}{1 + g _ {\varepsilon}}
\]

属于 I。对 \(x \notin \mathrm{C}_{\varepsilon}\)，有

\[
\left| f (x) - f _ {\varepsilon} (x) \right| \leqslant 2 \varepsilon ,
\]

而对 \(x\in \mathrm{C}_{\varepsilon}\)，有

\[
| f (x) - f _ {\varepsilon} (x) | = \frac {| f (x) |}{1 + g _ {\varepsilon} (x)} \leqslant \varepsilon | f (x) |.
\]

于是 \(f_{\varepsilon}\) 当 \(\varepsilon\) 趋于 0 时在 X 上一致收敛于 f。因此我们有 \(f \in \bar{I}\)，而由于 I 是闭的，就有 \(f \in I\)。

推论 1. --- 对每个 \(x \in X\)，用 \(I_x\) 表示由 \(f \in \mathcal{C}_0(X)\)（在 \(x\) 处为 0）所成的集合。则 \(x \mapsto I_x\) 是从 \(X\) 到 \(\mathcal{C}_0(X)\) 的极大闭理想所成的集合上的一个双射。这些理想都是正则的。

这立即由命题 1 得出。

用 X' 表示 X 的 Alexandroff 紧化，即由 X 添加一个无穷远点 \(\omega_{X}\) 而得到的紧空间 (TG, I, pp. 67–68)。代数 \(\mathcal{C}_{0}(X)\) 与 X' 上在 \(\omega_{X}\) 处为 0 的复值连续函数所成的 Banach 代数视为同一。

对每个 \(x \in X'\)，我们用 \(ev_{x}\) 表示 \(\mathcal{C}_{0}(X)\) 的由 \(\mathrm{ev}_{x}(f) = f(x)\)（对一切 \(f \in \mathcal{C}_{0}(X)\)）所定义的特征标。

推论 2. --- 映射 \(x \mapsto \mathrm{ev}_x\) 是从 \(\mathbf{X}'\) 到 \(\mathsf{X}'(\mathcal{C}_0(\mathbf{X}))\) 上的一个同胚，而它在 \(\mathbf{X}\) 上的限制是从 \(\mathbf{X}\) 到 \(\mathsf{X}(\mathcal{C}_0(\mathbf{X}))\) 上的一个同胚。此外，弱拓扑与 Jacobson 拓扑在 \(\mathsf{X}(\mathcal{C}_0(\mathbf{X}))\) 上重合。

映射 ev：\(x \mapsto ev_{x}\)（从 \(X'\) 到 \(\mathsf{X}'(\mathcal{C}_{0}(\mathbf{X}))\) 中的映射）是单射。由 I, p.~30 的推论 1 和定理 2，它是满射。它是连续的，因为对每个函数 \(f \in \mathcal{C}_{0}(\mathbf{X})\) 及 R 的每个开集 U，我们有

\[
\mathrm{ev} ^ {- 1} (\{\chi \in \mathsf {X} ^ {\prime} (\mathcal {C} _ {0} (\mathsf {X})) \mid \chi (f) \in \mathrm{U} \}) = f ^ {- 1} (\mathrm{U})
\]

它在 X 中是开的。由于 \(X'\) 是紧的，映射 ev 因此是一个同胚。于是 ev 在 X 上的限制是一个到 \(\mathsf{X}(\mathcal{C}_{0}(\mathsf{X}))\) 上的同胚。

若 F 是 \(\mathsf{X}(\mathcal{C}_{0}(\mathbf{X}))\) 的一个弱闭子集，则它经由同胚 ev 对应于 X 的一个闭子集 \(\Phi\)；确切地说，由命题 1，我们有 \(\mathrm{F} = \{\chi \in \mathsf{X}(\mathcal{C}_{0}(\mathbf{X})) \mid \mathrm{I}_{\Phi} \subset \mathrm{Ker} \, \chi\}\)，它对于 Jacobson 拓扑是闭的。

推论 3. --- 假定 X 是紧的。则映射 \(x \mapsto \mathrm{ev}_x\) 是从 X 到 \(\mathsf{X}(\mathcal{C}(\mathbb{X}))\) 上的一个同胚。弱拓扑与 Jacobson 拓扑在 \(\mathsf{X}(\mathcal{C}(\mathbb{X}))\) 上重合。
% semantic-fragments: legacy-input-seam
\relax{}
